\newcommand{\maiheader}{
\begin{center}
\textbf{МИНИСТЕРСТВО НАУКИ И ВЫСШЕГО ОБРАЗОВАНИЯ}\\
\textbf{РОССИЙСКОЙ ФЕДЕРАЦИИ}\\[0.4em]
\textbf{ФЕДЕРАЛЬНОЕ ГОСУДАРСТВЕННОЕ БЮДЖЕТНОЕ ОБРАЗОВАТЕЛЬНОЕ}\\
\textbf{УЧРЕЖДЕНИЕ ВЫСШЕГО ОБРАЗОВАНИЯ}\\
\textbf{«МОСКОВСКИЙ АВИАЦИОННЫЙ ИНСТИТУТ}\\
\textbf{(национальный исследовательский университет)»}
\end{center}
}

\newcommand{\fieldline}[2]{\underline{\makebox[#1][l]{#2}}}

\thispagestyle{empty}
\small
\maiheader

\vspace{1.0em}
\noindent\textbf{Институт (Филиал) №8 «Компьютерные науки и прикладная математика»}\\
\textbf{Кафедра 806}\\
\textbf{Группа М8О-403Б-22}\\
\textbf{Направление подготовки 01.03.02 Прикладная математика и информатика}\\[0.3em]
\textbf{Профиль Информатика}\\[0.3em]
\textbf{Квалификация бакалавр}\\
\rule{2.5cm}{0.4pt}

\vspace{1.4em}
\begin{flushright}
\begin{tabular}{@{}lll@{}}
\textbf{УТВЕРЖДАЮ} & & \\
Заведующий кафедрой №806 & \fieldline{2.4cm}{} & С.\,С. Крылов \\
& {\scriptsize (№ каф.) \hspace{2.0cm} (подпись) \hspace{1.0cm} (инициалы, фамилия)} & \\
\fieldline{0.8cm}{} \hspace{0.3em} \fieldline{2.0cm}{} 2026 г. & &
\end{tabular}
\end{flushright}

\begin{center}
\textbf{ЗАДАНИЕ}\\
\textbf{на выпускную квалификационную работу}\\
\textbf{бакалавра}
\end{center}

\vspace{0.6em}
\noindent
\begin{tabularx}{\textwidth}{@{}lX@{}}
\textbf{Обучающийся} & \fieldline{10.5cm}{Мудров Павел Федорович} \\
& {\scriptsize (фамилия, имя, отчество полностью)} \\
\textbf{Руководитель} & \fieldline{10.5cm}{Ляпина Светлана Юрьевна} \\
& {\scriptsize (фамилия, имя, отчество полностью)} \\
\multicolumn{2}{@{}l@{}}{\fieldline{15.2cm}{профессор}} \\
\multicolumn{2}{@{}l@{}}{{\scriptsize ученая степень, ученое звание, должность и место работы}}
\end{tabularx}

\vspace{0.8em}
\noindent\textbf{1. Наименование темы}\hspace{0.8em}
«Применение методов обучения с подкреплением (Reinforcement Learning, RL) для прогнозирования и оптимизации биржевых котировок»

\vspace{0.8em}
\noindent\textbf{2. Срок сдачи обучающимся законченной работы}\hspace{0.8em}\fieldline{2.3cm}{15.05.2026}

\vspace{0.8em}
\noindent\textbf{3. Задание и исходные данные к работе}\\
Разработать единый экспериментальный стенд сравнения baseline-стратегий, статистических подходов, моделей машинного обучения и агентов обучения с подкреплением на биржевых котировках. Реализовать программный комплекс на Python с загрузкой данных, формированием признаков, обучением моделей, расчетом метрик и визуализацией результатов.

\vspace{0.3em}
\noindent\textbf{Перечень иллюстративно-графических материалов *при наличии:}

\vspace{0.3em}
\begin{center}
\begin{tabularx}{0.92\textwidth}{|c|X|c|}
\hline
№ п/п & Наименование & Количество листов \\
\hline
\rule{0pt}{1.4em} & & \\
\hline
\rule{0pt}{1.4em} & & \\
\hline
\rule{0pt}{1.4em} & & \\
\hline
\end{tabularx}
\end{center}

\newpage
\thispagestyle{empty}
\small

\noindent\textbf{4. Перечень подлежащих разработке разделов и этапы выполнения работы}

\vspace{0.4em}
\begin{center}
\setlength{\tabcolsep}{3pt}
\begin{tabularx}{0.92\textwidth}{|c|X|>{\centering\arraybackslash}p{2.3cm}|>{\centering\arraybackslash}p{2.6cm}|p{2.7cm}|}
\hline
№ п/п & Наименование раздела или этапа & Трудоемкость в \% от полной трудоемкости квалификационной работы & Срок выполнения & Примечание \\
\hline
1 & Анализ предметной области, аналогов и постановка исследовательской проблемы & 15 & 15.02.2026 & Выполнено на основе анализа литературы \\
\hline
2 & Проектирование экспериментального стенда и постановка задач ML и RL & 20 & 01.03.2026 & Согласовано с руководителем \\
\hline
3 & Реализация загрузки данных, генерации признаков, baseline-стратегий и ML-моделей & 25 & 20.03.2026 & Реализовано на языке Python \\
\hline
4 & Реализация RL-среды, обучение DQN и PPO, экспериментальное моделирование & 25 & 10.04.2026 & Проведено тестирование \\
\hline
5 & Подготовка выводов, оформление текста ВКР, рисунков и приложений & 15 & 15.05.2026 & Итоговое оформление работы \\
\hline
\end{tabularx}
\end{center}

\vspace{0.3em}
\noindent\textbf{5. Исходные материалы и пособия}
\begin{enumerate}[leftmargin=1.3cm,itemsep=0pt,topsep=0pt,label=\arabic*.]
\item Sutton R.\,S., Barto A.\,G. \textit{Reinforcement Learning: An Introduction}.
\item Hyndman R.\,J., Athanasopoulos G. \textit{Forecasting: Principles and Practice}.
\item Mnih V. et al. \textit{Human-level Control through Deep Reinforcement Learning}.
\item Schulman J. et al. \textit{Proximal Policy Optimization Algorithms}.
\item Документация \textit{stable-baselines3}, \textit{Gymnasium}, \textit{pandas}, \textit{scikit-learn}, \textit{yfinance}.
\end{enumerate}

\vspace{0.2em}
\noindent\textbf{6. Дата выдачи задания}\hspace{0.8em}\fieldline{2.3cm}{11.02.2026}

\vspace{0.6em}
\begin{tabularx}{\textwidth}{@{}X X@{}}
Руководитель \fieldline{4.3cm}{} & Задание принял к исполнению \fieldline{4.3cm}{} \\
{\scriptsize (подпись)} & {\scriptsize (подпись)}
\end{tabularx}

\newpage
\thispagestyle{empty}
\small
\maiheader

\vspace{1.0em}
\noindent\textbf{Институт (Филиал) №8 «Компьютерные науки и прикладная математика»}\\
\textbf{Кафедра 806}\\
\textbf{Группа М8О-403Б-22}\\
\textbf{Направление подготовки 01.03.02 Прикладная математика и информатика}\\[0.3em]
\textbf{Профиль Информатика}\\[0.3em]
\textbf{Квалификация бакалавр}

\vspace{1.8em}
\begin{center}
\textbf{ВЫПУСКНАЯ КВАЛИФИКАЦИОННАЯ РАБОТА}\\
\textbf{БАКАЛАВРА}
\end{center}

\vspace{0.8em}
\noindent
На тему: «Применение методов обучения с подкреплением (Reinforcement Learning, RL) для прогнозирования и оптимизации биржевых котировок»

\vspace{2.2em}
\noindent
Автор ВКРБ \fieldline{4.0cm}{} Мудров Павел Федорович\\
\hspace*{3.6cm}(\fieldline{1.5cm}{}) \hspace{3.0cm}{\scriptsize (фамилия, имя, отчество полностью)}\\[0.7em]
Руководитель \fieldline{3.2cm}{} Ляпина Светлана Юрьевна\\
\hspace*{3.7cm}(\fieldline{1.5cm}{}) \hspace{2.9cm}{\scriptsize (фамилия, имя, отчество полностью)}\\[0.7em]
Консультант \fieldline{3.1cm}{} \fieldline{6.2cm}{}\\
\hspace*{3.7cm}(\fieldline{1.5cm}{}) \hspace{2.9cm}{\scriptsize (фамилия, имя, отчество полностью)}\\[0.7em]
Консультант \fieldline{3.1cm}{} \fieldline{6.2cm}{}\\
\hspace*{3.7cm}(\fieldline{1.5cm}{}) \hspace{2.9cm}{\scriptsize (фамилия, имя, отчество полностью)}\\[0.7em]
Рецензент \fieldline{3.4cm}{} \fieldline{6.2cm}{}\\
\hspace*{3.7cm}(\fieldline{1.5cm}{}) \hspace{2.9cm}{\scriptsize (фамилия, имя, отчество полностью)}

\vspace{0.8em}
\noindent
\textbf{К защите допустить}\\
Заведующий кафедрой № 806 \hspace{0.8em} Крылов Сергей Сергеевич \hspace{0.8em} (\fieldline{1.6cm}{})\\
\hspace*{5.0cm}{\scriptsize (№ каф.) \hspace{3.2cm} (фамилия, имя, отчество полностью)}\\
\fieldline{0.8cm}{} \hspace{0.3em} \fieldline{2.3cm}{} 2026 г.

\vspace{1.0em}
\begin{center}
Москва 2026
\end{center}

\newpage
\thispagestyle{empty}
\normalsize
\begin{center}
\textbf{СПИСОК ИСПОЛНИТЕЛЕЙ}
\end{center}

\noindent
В выполнении выпускной квалификационной работы принимали участие:

\vspace{0.8em}
\noindent
Обучающийся --- Мудров Павел Федорович

\noindent
Руководитель --- Ляпина Светлана Юрьевна, профессор

\noindent
Консультант --- TODO: заполнить сведения при наличии

\noindent
Рецензент --- TODO: заполнить сведения при наличии

\newpage
\thispagestyle{empty}
\normalsize
\begin{center}
\textbf{РЕФЕРАТ}
\end{center}

\noindent
Выпускная квалификационная работа бакалавра состоит из 57 страниц, 28 использованных источников, 7 таблиц, 5 рисунков и одного приложения.

\vspace{0.5em}
\noindent
\textbf{ОБУЧЕНИЕ С ПОДКРЕПЛЕНИЕМ, БИРЖЕВЫЕ КОТИРОВКИ, АЛГОРИТМИЧЕСКАЯ ТОРГОВЛЯ, DQN, PPO, BASELINE-СТРАТЕГИИ, МАШИННОЕ ОБУЧЕНИЕ, ФИНАНСОВЫЕ МЕТРИКИ, EQUITY CURVE, MAX DRAWDOWN}

\vspace{0.5em}
\noindent
Работа посвящена исследованию применения методов обучения с подкреплением для прогнозирования и оптимизации торговых решений на биржевых котировках. В отличие от подходов, ориентированных исключительно на прогноз будущей цены, в работе рассматривается задача последовательного выбора действий \textit{buy}, \textit{sell} и \textit{hold} с учетом состояния рынка и динамики портфеля.

\noindent
В теоретической части рассмотрены особенности финансовых временных рядов, классические статистические методы, модели машинного обучения, нейросетевые подходы и современные RL-алгоритмы. Особое внимание уделено сравнению аналогов и формулированию проблемы, состоящей в том, что высокая точность прогноза не гарантирует практической прибыльности стратегии. На этой основе обоснована необходимость построения единого экспериментального стенда, в котором baseline-стратегии, ML-модели и RL-агенты оцениваются на одних и тех же данных и по единому набору метрик.

\noindent
В практической части описан программный проект на Python, включающий загрузку данных через \textit{yfinance}, генерацию признаков на основе OHLCV-рядов, реализацию стратегий \textit{Buy \& Hold} и \textit{Moving Average Crossover}, обучение моделей \textit{Logistic Regression}, \textit{Random Forest}, \textit{Gradient Boosting}, а также обучение агентов \textit{DQN} и \textit{PPO} в среде \textit{Gymnasium}. Для оценки результатов используются как ML-метрики, так и финансовые показатели: доходность, волатильность, коэффициент Шарпа, максимальная просадка и устойчивость стратегии.

\noindent
Результаты экспериментального моделирования показывают, что RL-подходы целесообразно рассматривать как инструмент оптимизации последовательных торговых решений, однако их качество существенно зависит от функции вознаграждения, набора признаков и рыночного режима. В условиях исследуемого стенда PPO-агент обеспечивает более устойчивый риск-скорректированный результат по сравнению с рядом альтернатив, тогда как стратегия \textit{Buy \& Hold} сохраняет преимущество по абсолютной доходности на выраженном восходящем рынке.
